\scanpage{16}
Une telle $V$ a donc une base canonique (commençant par le $a_0$ donné)
\[
(a_0,a_1,\,-\,,a_n)
\]
de telle façon que \inserted{le} $a_i$ \unsure{fasse|forme} base de $L_i$ \inserted{$0\leq i\leq n$}, \deleted{\unsure{d'où}} qu'on a
\[
\tau_i=\mathrm{id}+\deleted{\unsure{a}}a_i^\vee\otimes a_i\qquad (
\]
avec $a_i^\vee\in\Gamma V^\vee$ bien déterminé, et
\[
\langle a_i,a_{i+1}^\vee\rangle=1\qquad 0\leq i\leq n-1
\]
(On a alors exprimé déjà \deleted{1), 2)} sur $(\tau_i^V)$ 1), 2) 3), il \deleted{On aura alors \unsure{unic.|deux}} reste à exprimer \deleted{\unsure{3) et}} la commutation des $\tau_i^V,\tau_j^V$ pour $i,j$ non consécutifs, ce qui s'exprime par
\[
L_i\subset\unsure{\mathrm{Hyp}|\mathrm{Ker}}(a'_j)\qquad\text{si}\quad 0\leq i,j\leq n,\quad i,j\text{ non consécutifs.}
\]

On aura donc
\[
\left\{\begin{aligned}
 a_0^\vee&=\deleted{\unsure{(\lambda_0-1)a_0^*+\beta_0a_1^*}}(\lambda_0-1)a_0^*+\beta_0a_1^*\\
 a_i^\vee&=\deleted{\unsure{\beta|\xi}}a_{i-1}^*+(\lambda_i-1)a_i^*+\beta_i a_{i+1}^* & \unsure{0|1}\unsure{<|\leq}i\leq n-1\\
 a_n^\vee&=a_{n-1}^*+(\lambda_n-1)a_n^*
\end{aligned}\right.
\]
\margin{\uline{NB} $a_0^\vee=0$ sss $\lambda_0=1,\beta_0=0$. Donc $\tau_0$ ps. refl. stricte sss \deleted{\unsure{$\lambda_0\ne1$}} \inserted{en tt} $s\in S$, on a $\lambda_0(s)\ne1$ \uline{ou} $\beta_0(s)\ne0$, i.e. sss $(\lambda_0-1)\mathcal O_S+\beta_0\mathcal O_S=\mathcal O_S$}

\uline{Scholie} On a un foncteur $C\longrightarrow C'_0$ évident,
\[
(E,\tau_*,s_0)\longmapsto(V\deleted{\unsure{,|_{*}}}
\]
qui est une équivalence de catégories

$\boxed{C''_0}$ \deleted{\unsure{Corollaire}} Considérons la catégorie des \deleted{\unsure{couples}} \deleted{couples} triples \margin{$n\geq2$} $(C,\varphi_G^C,h_0)$, où $C$ est un fibré projectif de rg $n$, $\varphi_{\unsure{G|g}}^{C}\colon{C}\longrightarrow\mathrm{Aut}(C)$ hom de groupe, $h_0$ section de $C$, avec conditions

1) $\forall\,0\leq i\leq n$, $\tau_i^C\overset{\mathrm{dfn}}=\varphi_G(\underline{\tau}_i)$ est une pseudo-réflexion, qui est stricte si $1\leq i\leq n$, et associe un centre $h_0$ si $i=0$ (\deleted{\unsure{$\tau_0$}} pas néc. stricte) -- d'où $h_i,H_i$ $(1\leq i\leq n)$

2) Les $h_i$ $0\leq i\leq n$ lin. indép. en chaque fibre

3) $h_i\notin H_{i+1}$ \deleted{\unsure{en}} tte fibre $(0\leq i\leq n-1)$
% archive sheet number: 15

